\documentclass[10pt,a4paper]{amsart}
\usepackage[margin=19mm]{geometry}
\usepackage{amsmath,amssymb,mathtools}
\usepackage[scr=rsfs,cal=euler]{mathalfa}
\usepackage{xfrac}
\usepackage[unicode,hidelinks,bookmarks=false]{hyperref}
\newcommand{\ie}{i.e.\ }
\newcommand{\cf}{cf.\ }
\newcommand{\resp}{respectively\ }
\newcommand{\Subsection}{\subsection}
\providecommand{\nobreakdash}{\nobreak\hbox{-}}
\setlength{\emergencystretch}{2em}
\multlinegap0pt
\usepackage{bm}
\usepackage{bbm}
\usepackage{dsfont}
\usepackage{xspace}
\usepackage{cancel}
\usepackage{mathtools}
\usepackage{amsthm}
\makeatletter
\@ifundefined{teo}{\newtheorem{teo}{Teo}}{}
\@ifundefined{lem}{\newtheorem{lem}[teo]{Lemma}}{}
\makeatother
\def\de{{\rm d}}
\newcommand{\R}{\bbR}
\newcommand{\bbN}{{\mathbb{N}}}
\newcommand{\bbR}{{\mathbb{R}}}
\newcommand{\hol}[1]{\"{#1}}
\newcommand{\supp}{\mathrm{supp}}
\title[Minimization of Degenerate Nonlinear Functionals under Radia]{Minimization of Degenerate Nonlinear Functionals under Radial Symmetry}
\author{Valeria Chiad\`o Piat, Virginia De Cicco, Anderson Melchor Hernandez}
\date{Source: arXiv:2507.20603v1 (CC BY 4.0)}
\begin{document}
\maketitle

\noindent\emph{Source and licence.} Minimization of Degenerate Nonlinear Functionals under Radial Symmetry. Version arXiv:2507.20603v1 (CC BY 4.0), distributed under the Creative Commons Attribution 4.0 International licence, as recorded in the arXiv licence metadata for this exact version and shown on the abstract page https://arxiv.org/abs/2507.20603v1. Full rights notice: licenses/arxiv-2507.20603-CC-BY-4.0.txt.\par\smallskip

\noindent\textit{Editorial scope.} Counted unit: the Lem whose source label is \texttt{lemma2} in the article (source0.tex lines 447--466, source label \texttt{lemma2}), delivered together with the article's own complete original proof of it (source0.tex lines 467--559, one \texttt{proof} environment of 93 lines). No mathematics is written, repaired or re-derived: statement and proof are the article's own text line for line. Required context retained uncounted: decomposet (lines 270--274); newset (lines 292--295). Every definition, display and label of those lines is retained unchanged. Omitted from this package and not redistributed: 1--269, 275--291, 296--446, 560--1263. Nothing inside a retained excerpt is omitted or altered: every retained body is delivered exactly as the article's lines are written, so no omission receipt of any kind accompanies this package. Citations: the retained excerpts carry no citation command at all, so no bibliography entry of the article is transcribed and no reference is redistributed. Reference adaptation: none was needed and none was made; every reference inside the delivered unit resolves inside it, so no unresolved reference or citation is produced. Printed slips kept literal: no printed word, symbol, hypothesis or number of the counted statement or of its proof was altered. Layout and class: the article's own class is \texttt{amsart} with packages including amsmath, amsfonts, mathtools, amssymb, mathrsfs, subcaption, bookmark, geometry, epstopdf, amscd, xcolor, graphicx, yfonts, esint; the wrapper is the lane's pinned \texttt{amsart} editorial wrapper at 10pt on A4, which restates the theorem-like environments the retained text needs and adds the article's own macro definitions that the retained text needs (\texttt{\textbackslash R}, \texttt{\textbackslash bbN}, \texttt{\textbackslash bbR}, \texttt{\textbackslash de}, \texttt{\textbackslash hol}, \texttt{\textbackslash supp}), with extra wrapper packages (bm, bbm, dsfont, xspace, cancel, mathtools). The delivered numbering is the unit's own. Assets: no retained span contains a figure environment, a graphics-inclusion command, a PDF-page inclusion, a table or a TikZ picture; no image file is copied into this package, the package declares no asset dependency and no image file is redistributed with it. Licence: version arXiv:2507.20603v1 of this article is distributed under the Creative Commons Attribution 4.0 International licence, as recorded in the arXiv licence metadata for this exact version and shown on the abstract page https://arxiv.org/abs/2507.20603v1. A full rights notice is delivered at licenses/arxiv-2507.20603-CC-BY-4.0.txt. Characters: every retained excerpt is pure ASCII, so the delivered engine, XeLaTeX with its default Latin Modern text font, reports no missing character.
\begin{align}
\label{decomposet}
I_{p,{\mathrm supp}(\eta)}:=\Big\{r\in({\mathrm supp}(\eta))^{\circ}:\,&\exists\,\epsilon >0\text{ such that }
\eta^{-\frac{1}{p-1}}\in L^1\left((r-\epsilon ,r+\epsilon )\right)\Big\},
\end{align}

\begin{align}
\label{newset}
I_{\alpha_{i},\beta_{i}}\coloneqq \{x\in I_{\Omega,w}: \alpha_{i}<\vert x\vert< \beta_{i}\}, \hskip 0,2cm\, a_{i}<\alpha_{i}<\beta_{i}\leq b_{i}, \hskip 0,2cm i=1,\ldots, N_{\eta}.
\end{align}

\begin{lem}[Fundamental convergence]\label {lemma2} 
Suppose that $\eta$ is finitely degenerate. Let $(u_h)_h\subset W^{1,1}(\Omega)$ be a sequence of radial functions such that
\begin{enumerate}
\item[(a)] $\displaystyle\sup_{h\in\bbN}\int_{I_{\Omega,w}}|\nabla u_h|^{p}w\,\de x<+\infty$\,,
\item[(b)] for every $i=\,1,\dots,N_\eta$ there exists $c_i$ such that $a_i< c_i<b_i$ and there exist finite the following limits
$$
\lim_{h\to+\infty}u_h(x)=d_i\in \bbR\, \hskip 0,1cm \text{for all $x\in I_{\Omega,w}$ such that $\vert x\vert=c_{i}$}.
$$
\end{enumerate}
Then there exists a subsequence $(u_{h_k})$ and a radial function 
$u:\,I_{\Omega,w}\to\bbR$ such that 
\begin{enumerate}
\item[(i)] $\displaystyle\lim_{k\to+\infty}u_{h_k}(x) =u(x) $ for every $x\in I_{\Omega,w}$\,,
\item[(ii)] $u\in {\rm{Dom}}_{{\mathrm r},w}$,
\item[(iii)]
\begin{align*}
\displaystyle\int_{I_{\Omega,w}}\vert \nabla u \vert^{p}w\de x\leq \liminf_{h_{k}\rightarrow +\infty}\displaystyle\int_{I_{\Omega,w}}\vert \nabla u_{h_{k}}\vert^{p}w\de x.
\end{align*}
\end{enumerate}
\end{lem}

\begin{proof}
 Let us note that, by assumption (b), $I_{\Omega,w}\neq\emptyset$. By hypothesis, we have that $(u_{h})_{h}$ is a sequence of radial functions. That is,
\begin{align*}
 u_{h}(x)=v_{h}(\vert x\vert), \hskip 0,2cm \text{for all $h\in\mathbb{N}$, $x\in \Omega$.}
 \end{align*}
 Notice that 
\begin{align*}
\frac{\partial u_{h}(x)}{\partial x_{j}}=v_{h}'(\vert x\vert)\frac{x_{j}}{r}, \hskip 0,2cm\, r=\vert x \vert,\, j=1,\ldots,d,
\end{align*}
 and thus 
\begin{align*}
\int_{I_{\Omega,w}}|\nabla u_h|^{p}w\,\de x&=\int_{I_{\Omega,w}}\left[\left\{\sum_{j=1}^{d}|v_{h}'(\vert x\vert)|^{2}\frac{\vert x_{j}\vert^{2}}{\vert x\vert^{2}}\right\}^{\frac{1}{2}}\right]^{p}w\,\de x\\
&=\int_{I_{\Omega,w}}|v_{h}'(\vert x\vert)|^{p}w\,\de x.
\end{align*}
Moreover, since $w$ is radial, we have that by the change of variable theorem
\begin{align*}
\displaystyle\int_{I_{\Omega,w}}|v_{h}'(\vert x\vert)|^{p}w\,\de x&=\displaystyle\int_{I_{\Omega,w}}|v_{h}'(\vert x\vert)|^{p}\eta(\vert x\vert)\,\de x\\
&=\omega_{d}\displaystyle\int_{I_{p,\supp(\eta)}}\hskip -0,5cm|v_{h}'(r)|^{p}r^{d-1}\eta(r)\,\de r,
\end{align*}
where 
\begin{align}
\label{misurasfera}
\omega_{d}\coloneqq \mathcal{H}^{d-1}(\mathbb{S}^{d}).
\end{align}
By (a), notice that 
\begin{align*}
\sup_{h\in \mathbb{N}}\displaystyle\int_{I_{p,\supp(\eta)}}\hskip -0,5cm|v_{h}'(r)|^{p}r^{d-1}\eta(r)\,\de r<+\infty.
\end{align*}
Then there exist a subsequence $(v_{h_k})_k$ of $(v_h)_h$, 
and  a function $v\in L^{p}(I_{p,\supp(\eta)},\eta)$ such that 
\begin{equation}\label{(2)}
v_{h_k}'\to v {\ \ \rm weakly \ in \ }L^{p}(I_{p,\supp(\eta)},\eta)\text{ as }k\to\infty\,,
\end{equation} 
\begin{equation*}
\frac{\partial u_{h_{k}}}{\partial x_{j}}\to v(\vert x\vert)\frac{x_{j}}{\vert x\vert}{\ \ \rm weakly \ in \ }L^{p}(I_{\Omega,w},w)\text{ as }k\to\infty\,, j=1,\ldots,d.
\end{equation*} 
Moreover, let us observe that
\begin{equation*}\label{(3)}
L^{p}_{\rm{loc}}(I_{\Omega,w},w)\subset L^1_{\rm{loc}}(I_{\Omega,w})\,, \hskip 0,5cm L^{p}_{\rm{loc}}(I_{p,\supp(\eta)},\eta)\subset L^1_{\rm{loc}}(I_{p,\supp(\eta)}).
\end{equation*}
Indeed, let us notice that by the change of variable
\begin{align*}
\displaystyle\int_{I_{\Omega,w}} w^{-\frac{p'}{p}} \de x=\omega_{d}\displaystyle\int_{I_{p,\supp(\eta)}} \eta^{-\frac{p'}{p}}r^{d-1} \de r,
\end{align*}
where $\frac{1}{p}+\frac{1}{p'}=1$ . Then for $i=1,\ldots, N_{\eta}$, and for each $K\Subset (a_{i},b_{i})$, we get by \eqref{decomposet}, and since $\frac{p'}{p}=\frac{1}{p-1}$ that
\begin{align*}
\displaystyle\int_{K} \eta^{-\frac{p'}{p}}r^{d-1} \de r\leq b_{i}^{d-1}\displaystyle\int_{K} \eta^{-\frac{p'}{p}}\de r<+\infty.
\end{align*}
Let us take a compact set $\mathcal{K}\Subset I_{\alpha_{i},\beta_{i}}$, $i=1,\ldots, N_{\eta}$ with $I_{\alpha_{i},\beta_{i}}$ as defined in \eqref{newset}. Then by H\hol{o}lder's inequality
\begin{align}
\label{usefulbound}
\displaystyle\int_{\mathcal{K}}\vert  z\vert \de x\leq \left(\displaystyle\int_{\mathcal{K}}\vert  z\vert^{p}w\de x\right)^{\frac{1}{p}}\left(\displaystyle\int_{\mathcal{K}} w^{-\frac{p'}{p}} \de x\right)^{\frac{1}{p'}}<+\infty
\end{align}
for any $z\in L^{p}(\mathcal{K},w)$. Notice that from \eqref{(2)} and \eqref{usefulbound}, we get that $v\in L^1_{\rm{loc}}(I_{\alpha,\beta};\mathbb{R})$ and
\begin{equation*}\label{4}
\displaystyle\int_{I_{\alpha,\beta}} \left\vert \frac{\partial u_{h_k}}{\partial x_{j}}- v(\vert x\vert)\frac{x_{j}}{\vert x\vert}\right\vert \,\de x\rightarrow 0\, \text{ as }k\to\infty\,,
\end{equation*}
for each $[\alpha,\beta]\subset\,I_{p,{\mathrm supp}(\eta)}$.
Let us consider $u:\,\Omega\to\R$ defined in the following way:  let $i=\,1,\dots,N_\eta$, and consider an interval $(a_{i}, b_{i})$ and define $I_{a_{i},b_{i}}$ as in \eqref{newset}, so that we let
\begin{equation*}
u^{i}(x):=
\begin{cases}
\ \ \ \ \ 0 \ \ \ \ &\text{ if } x\in\Omega\setminus I_{a_{i},b_{i}},\,
\\
\displaystyle d_{i}+\int_{c_{i}}^{\vert x\vert}v(r)\de r&\text{ if } c_{i}\leq \vert x\vert<b_{i},\\
\displaystyle d_{i}+\int_{\vert x\vert}^{c_{i}}v(r)\de r&\text{ if } a_{i}\leq \vert x\vert<c_{i}.
\end{cases}
\end{equation*}
Then we define
$$
u(x)=\sum_{i=1}^{N_{\eta}} u^i(x)  \chi_{I_{a_{i},b_{i}}}(x) .
$$
In what follows, we aim to prove that
\begin{equation*}\label{aimed}
u\in W^{1,1}_{\text {loc}}(I_{\Omega, w}).
\end{equation*}
As before, let us take some $\mathcal{K}\Subset I_{\alpha_{i},\beta_{i}}$, $i=i,\ldots, N_{\eta}$. First, notice that by the fundamental theorem of calculus, one has that
\begin{align*}
\frac{\partial u}{\partial x_{j}}=v(\vert x\vert)\frac{x_{j}}{\vert x\vert},
\end{align*}
and then
\begin{align*}
\sum_{j=1}^{d}\displaystyle\int_{\mathcal{K}}\left\vert \frac{\partial u}{\partial x_{j}}\right\vert\de x\leq\displaystyle\int_{\mathcal{K}} \vert v(\vert x \vert)\vert \de x
&\leq \omega_{d}\displaystyle\int_{a_{i}}^{b_{i}}\hskip -0,2cm r^{d-1}v(r)\de r<+\infty,
\end{align*}
where $\omega_{d}$ is defined according \eqref{misurasfera}. Moreover, let us notice that 
\begin{align}\label{finitamente1}
\displaystyle \int_{I_{\Omega,w}}\vert \nabla u\vert^{p}w\de x=\displaystyle \int_{I_{\Omega,w}}\vert v'(\vert x\vert)\vert^{p}w\de x&=\omega_{d}\int_{I_{p,\supp(\eta)}}\vert v'(r)\vert^{p} r^{d-1}\eta(r) \de r\\
&\leq \omega_{d}\liminf_{h_{k}}\displaystyle\int_{I_{p,\supp(\eta)}}\vert v_{h_{k}}'(r)\vert^{p} r^{d-1}\eta(r) \de r\nonumber\\
&=\liminf_{h_{k}}\displaystyle\int_{I_{\Omega,w}}\vert \nabla u_{h_{k}}\vert^{p}w\de x<+\infty,\nonumber
\end{align}
where in \eqref{finitamente1} we have used that $\eta$ is finitely degenerate, and in the last inequality we have used the radiality and assumption (a).
\end{proof}

\end{document}
